\section{An ordinary finite presentation for the Steinberg group}
\label{sec:finite-presentation}

Let $R$ be a unital associative algebra over $\F_2$ and let $n\geq3$.
Recall that $\St_n(R)$ is generated by the root symbols $x_{ij}(a)$
with relations \eqref{fp:steinberg-addition}--\eqref{fp:steinberg-product}.
We continue to use $[g,h]=ghg^{-1}h^{-1}$. In the commutation relation,
both generators have unequal row and column indices; in particular,
roots in a single position commute. The addition relation gives
$x_{ij}(0)=1$ and $x_{ij}(a)^2=1$.

Write $e_{ij}(a)$ for the matrix with diagonal entries $1$, entry $a$ in
position $(i,j)$, and all other entries $0$.  These matrices satisfy
Equations~\eqref{fp:steinberg-addition}--\eqref{fp:steinberg-product},
with $x$ replaced by $e$.  Let $\GL_n(R)$ be the group of invertible
$n\times n$ matrices over $R$, and let $\E_n(R)$ be the subgroup generated
by the matrices $e_{ij}(a)$.  We therefore have a surjective homomorphism
\begin{equation}
 \pi_n:\St_n(R)\longrightarrow\E_n(R),\qquad
 x_{ij}(a)\longmapsto e_{ij}(a).
 \label{fp:matrix-map}
\end{equation}
We write
\begin{equation}
 J_n(R)=\Ker\pi_n.
 \label{fp:kernel-definition}
\end{equation}
When the algebra is understood, we omit $R$ from these group names.

\subsection{The finite presentation}

Krsti\'c and McCool prove finite presentability of Steinberg groups
over finitely presented unital rings for $n\geq4$
\cite[Theorem~3]{KrsticMcCool}. The explicit characteristic-two
presentation below follows their method of extending monomial root
relations by word length and then imposing the algebra relators
\cite[Section~3]{KrsticMcCool}. We give the induction for $n\geq5$
and prove that the finite relations imply every Steinberg relation.
The algebra construction supplies the required input.

Let $\mathcal A$ be a finite alphabet.  Its free monoid
$\mathcal A^*$ includes the empty word $\varepsilon$, of length zero.
The free unital associative algebra
$T=\F_2\langle\mathcal A\rangle$ has $\mathcal A^*$ as a vector-space
basis, with $\varepsilon$ representing $1$ and multiplication given by
concatenation.  Thus each $f\in T$ has a unique expression
\[
 f=\sum_{w\in S_f}w
\]
for a finite set $S_f\subseteq\mathcal A^*$.  We fix an ordering of the
alphabet and use length followed by lexicographic order to order such
finite sets.  The order will only specify literal words in a finite group
presentation; the resulting root elements will commute.

\begin{theorem}[Finite presentation of the Steinberg group]
\label{thm:finite-presentation}
Let
\[
 R=T/\mathfrak a,
 \qquad
 T=\F_2\langle\mathcal A\rangle,
 \qquad
 \mathfrak a=(f_1,\ldots,f_s)
\]
be a finite presentation as a unital associative $\F_2$-algebra, where
$\mathfrak a$ denotes the two-sided ideal generated by the listed
polynomials.  If $n\geq5$, then $\St_n(R)$ has the following ordinary
finite group presentation.

The generators are
\[
 g_{ij}(w),\qquad i\neq j,\quad w\in\mathcal A^*,\quad |w|\leq2.
\]
To specify the relators, define a formal group word $\xi_{ij}(w)$ in
these generators for every $w\in\mathcal A^*$.  For $|w|\leq2$, put
$\xi_{ij}(w)=g_{ij}(w)$.  For $|w|>2$, write $w=av$, with
$a\in\mathcal A$ its first letter, choose the least
$h\in\{1,\ldots,n\}\setminus\{i,j\}$, and put
\[
 \xi_{ij}(w)=[\xi_{ih}(a),\xi_{hj}(v)].
\]
This recursion terminates and gives a finite group word for each $w$.
The defining relations are
\begin{align}
 g_{ij}(w)^2&=1,
 && |w|\leq2, \label{fp:finite-square}\\
 [g_{ij}(u),g_{kl}(v)]&=1,
 && |u|,|v|\leq2,\quad j\neq k,\ i\neq l,
 \label{fp:finite-commutation}\\
 [g_{ij}(u),g_{jk}(v)]&=g_{ik}(uv),
 && |u|+|v|\leq2,\quad i,j,k\text{ pairwise distinct},
 \label{fp:finite-product}\\
 \prod_{w\in S_{f_t}}\xi_{ij}(w)&=1,
 && 1\leq t\leq s,\quad i\neq j.
 \label{fp:finite-algebra-relators}
\end{align}
The product in Equation~\eqref{fp:finite-algebra-relators} is taken in
the fixed order above, and an empty product is $1$.  The isomorphism to
$\St_n(R)$ sends $g_{ij}(w)$ to $x_{ij}(\overline w)$, where
$\overline w$ is the image of $w$ in $R$.
\end{theorem}

The proof constructs a root for a longer monomial as a commutator of
roots for shorter pieces. Spare coordinates will show that the result
is independent of the split and the intermediate index. At each new
length we then establish commutation before treating products with the
empty word. The following group identity makes the split comparison;
its centrality conclusion is proved, rather than assumed.

\begin{lemma}[Three-factor commutator identity]
\label{fp:three-factor}
Let $X,Y,Z$ belong to a group, and put $A=[X,Y]$ and $B=[Y,Z]$.
Assume that $A$ commutes with $X$ and $Y$, that $B$ commutes with $Y$
and $Z$, and that $A$ commutes with $B$ and $X$ commutes with $Z$.
Then
\[
 [A,Z]=[X,B].
\]
Moreover, this common element centralizes each of $A,Z,X,B$.
\end{lemma}

\begin{proof}
Let $H_1=\langle A,Z\rangle$ and $H_2=\langle X,B\rangle$.
Every generator of $H_1$ commutes with every generator of $H_2$ by the
assumptions, so these two subgroups centralize one another.
Conjugation of $[X,Z]=1$ by $Y$, together with
\[
 YXY^{-1}=A^{-1}X,
 \qquad
 YZY^{-1}=BZ,
\]
gives $[A^{-1}X,BZ]=1$.  Since $B$ and $Z$ commute, and since the
factors from $H_1$ commute with those from $H_2$, expansion of this
commutator gives
\[
 1=[A^{-1}X,ZB]=[A^{-1},Z][X,B].
\]
Set $C=[A^{-1},Z]$ and $D=[X,B]$.  Then $C=D^{-1}$ belongs to both
$H_1$ and $H_2$.  An element in their intersection centralizes each
subgroup, so $C$ commutes with $A,Z,X,B$.  Finally, direct expansion
gives
\[
 [A,Z]^{-1}=A[A^{-1},Z]A^{-1}=ACA^{-1}=C=D^{-1}.
\]
This proves both assertions.
\end{proof}

\begin{proof}[Proof of Theorem~\ref{thm:finite-presentation}]
First omit Equation~\eqref{fp:finite-algebra-relators}, and denote the
resulting finitely presented group by $P_0$.  We will prove
$P_0\cong\St_n(T)$, and then impose the algebra relators.

\paragraph{Construction of all monomial roots.}
We construct elements $y_{ij}(w)\in P_0$ for every word $w$.  For
$|w|\leq2$, set $y_{ij}(w)=g_{ij}(w)$.  The precise induction
assertion at an integer $N\geq2$ is the following:
\begin{enumerate}
 \item The elements $y_{ij}(w)$ are defined for all $|w|\leq N$.
 For every positive-length split $w=uv$ with $|w|\leq N$ and every
 $h\notin\{i,j\}$, they satisfy
 \begin{equation}
  y_{ij}(w)=[y_{ih}(u),y_{hj}(v)].
  \label{fp:positive-split}
 \end{equation}
 In particular, the value is independent of both the split and $h$.
 \item Every $y_{ij}(w)$ with $|w|\leq N$ is an involution.
 \item Whenever $|u|,|v|\leq N$, $j\neq k$, and $i\neq l$, one has
 \begin{equation}
  [y_{ij}(u),y_{kl}(v)]=1.
  \label{fp:all-monomial-commutations}
 \end{equation}
 There is no restriction on $|u|+|v|$ beyond the consequence
 $|u|+|v|\leq2N$.
 \item Whenever $|u|+|v|\leq N$ and $i,j,k$ are pairwise distinct,
 \begin{equation}
  [y_{ij}(u),y_{jk}(v)]=y_{ik}(uv).
  \label{fp:all-monomial-products}
 \end{equation}
 This assertion includes $u=\varepsilon$ or $v=\varepsilon$.
\end{enumerate}
At $N=2$ these assertions follow directly from
Equations~\eqref{fp:finite-square}--\eqref{fp:finite-product}.

Suppose the assertions hold at $N-1$, where $N\geq3$, and fix a word
$w$ of length $N$ and distinct indices $i,j$.  Each positive split
$w=uv$ and each $h\notin\{i,j\}$ gives a candidate
$[y_{ih}(u),y_{hj}(v)]$, since both factors have already been defined.
We first show that all these candidates agree.

Write $w=uvz$ with all three pieces nonempty, and choose distinct
indices $i,a,b,j$.  Set
\[
 X=y_{ia}(u),\qquad Y=y_{ab}(v),\qquad Z=y_{bj}(z).
\]
The product relations at the preceding stage give
\[
 A=[X,Y]=y_{ib}(uv),\qquad
 B=[Y,Z]=y_{aj}(vz).
\]
Every word occurring in $X,Y,Z,A,B$ has length less than $N$.
The pairs $(A,X)$, $(A,Y)$, $(B,Y)$, $(B,Z)$, $(A,B)$, and
$(X,Z)$ all occupy the commuting positions in
Equation~\eqref{fp:all-monomial-commutations}.  Thus the induction
hypothesis verifies every assumption of Lemma~\ref{fp:three-factor}.
Its conclusion is
\begin{equation}
 [y_{ib}(uv),y_{bj}(z)]
 =
 [y_{ia}(u),y_{aj}(vz)].
 \label{fp:comparison}
\end{equation}
Only the two final brackets in this equation involve total word length
$N$.

For completeness, consider the graph whose vertices are pairs $(p,h)$,
where $1\leq p<N$ specifies a cut after the $p$th letter of $w$ and
$h\notin\{i,j\}$ specifies the intermediate index.  Equation~
\eqref{fp:comparison} equates the candidates at any two vertices whose
cut positions and intermediate indices are both different.  This graph
is connected.  Indeed, vertices with both coordinates different are
adjacent.  If $(p,a)$ and $(p,b)$ have the same first coordinate and
$a\neq b$, choose $q\neq p$ and an index $c$ different from $a,b$;
then
\[
 (p,a),\ (q,c),\ (p,b)
\]
is a path.  If $(p,a)$ and $(q,a)$ have the same second coordinate and
$p\neq q$, choose two distinct indices $b,c$, both different from $a$;
then
\[
 (p,a),\ (q,b),\ (p,c),\ (q,a)
\]
is a path.  Such choices exist because $N-1\geq2$ and $n-2\geq3$.
All candidates consequently agree.  Define $y_{ij}(w)$ to be their
common value.  This proves the positive-split assertion at length $N$,
for all choices of $w,i,j$, before making any other assertion at this
length.  It also shows, by induction, that $y_{ij}(w)$ is the element
represented by the explicitly specified word $\xi_{ij}(w)$.

\paragraph{Commuting pairs at the new length.}
We next prove Equation~\eqref{fp:all-monomial-commutations} for all
$|u|,|v|\leq N$.  Pairs with both lengths less than $N$ are already
covered.  For the other pairs, use induction on the total length
$|u|+|v|$, beginning at $N$ and ending at $2N$.  Commutation is
symmetric, so we may arrange that the first element in the desired
relation is $y_{ij}(u)$ with $|u|=N$; write the other element as
$y_{kl}(v)$.  The index assumptions are $j\neq k$ and $i\neq l$.
Choose
\[
 h\notin\{i,j,k,l\},
\]
which is possible since $n\geq5$, and split $u=u_1u_2$ into two
nonempty words.  The positive-split assertion gives
\[
 y_{ij}(u)=[y_{ih}(u_1),y_{hj}(u_2)].
\]
Both factors have commuting-pair positions with $y_{kl}(v)$:
for the first factor the needed inequalities are $h\neq k$ and
$i\neq l$, and for the second they are $j\neq k$ and $h\neq l$.
Their respective total lengths with $v$ are strictly smaller than
$|u|+|v|$.  Their commutations therefore follow either from the preceding
word-length stage or from the current induction on total length.
An element commuting with two elements also commutes with their
commutator.  This proves the desired relation and completes the
commutation induction, including pairs of two words of length $N$.

\paragraph{Involutions and empty-word inputs.}
Write a new root element as
\[
 C=y_{ij}(w)=[U,V],\qquad
 U=y_{ih}(u),\quad V=y_{hj}(v),
\]
using a positive split.  By the preceding word-length stage, $U$ and
$V$ are involutions.  The commutations just proved show that $C$
commutes with each of $U,V$, since their root positions are commuting
positions.  On the other hand, direct expansion using $U^2=V^2=1$
gives
\[
 UCU^{-1}=VUVU=C^{-1}.
\]
As $C$ commutes with $U$, it follows that $C=C^{-1}$ and $C^2=1$.

The product relation at total length $N$ is already established when
both inputs have positive length.  For a right empty-word input, fix
pairwise distinct $i,j,k$, write $w=uv$ with $u,v$ nonempty, and choose
$h\notin\{i,j,k\}$.  In Lemma~\ref{fp:three-factor}, take
\[
 X=y_{ih}(u),\qquad
 Y=y_{hj}(v),\qquad
 Z=y_{jk}(\varepsilon).
\]
Then
\[
 A=y_{ij}(w),\qquad B=y_{hk}(v).
\]
The formula for $A$ follows from positive-split independence; that for
$B$ uses the right empty-word relation at a smaller length.  All the
commutations in Lemma~\ref{fp:three-factor} are commuting-pair
relations among words of length at most $N$, and so have now been
proved.  The lemma and another positive split give
\[
 [y_{ij}(w),y_{jk}(\varepsilon)]
 =[y_{ih}(u),y_{hk}(v)]
 =y_{ik}(w).
\]
For a left empty-word input, apply the same lemma to
\[
 X=y_{ij}(\varepsilon),\qquad
 Y=y_{jh}(u),\qquad
 Z=y_{hk}(v).
\]
Here $A=y_{ih}(u)$ by the smaller-length left empty-word relation,
and $B=y_{jk}(w)$ by positive-split independence.  The lemma gives
\[
 [y_{ij}(\varepsilon),y_{jk}(w)]
 =[y_{ih}(u),y_{hk}(v)]
 =y_{ik}(w).
\]
This completes every part of the induction.  In particular, all
monomial Steinberg relations hold for words of arbitrary length.

\paragraph{From monomials to all algebra elements.}
For $f\in T$, define
\[
 y_{ij}(f)=\prod_{w\in S_f}y_{ij}(w).
\]
Same-position commutation makes this product independent of its order.
The involution relations cancel terms appearing twice, so
\[
 y_{ij}(f+g)=y_{ij}(f)y_{ij}(g).
\]
Likewise, the monomial commutations imply
$[y_{ij}(f),y_{kl}(g)]=1$ whenever $j\neq k$ and $i\neq l$.

For pairwise distinct $i,j,k$, each commutator of a monomial factor in
$y_{ij}(f)$ with a monomial factor in $y_{jk}(g)$ is
$y_{ik}(uv)$.  Every such resulting factor commutes with all factors of
both input products: positions $(i,k)$ commute with both $(i,j)$ and
$(j,k)$.  Expanding the commutator of the two finite products therefore
gives
\[
 [y_{ij}(f),y_{jk}(g)]
 =\prod_{u\in S_f}\prod_{v\in S_g}y_{ik}(uv)
 =y_{ik}(fg).
\]
For the first equality, the usual identities for the commutator of a
product reduce to multiplication of the individual commutators because
all those commutators centralize both input subgroups.  For the second,
same-position commutation permits rearrangement, and involutivity
implements exactly the cancellation of equal monomials over $\F_2$.
Thus the elements $y_{ij}(f)$ satisfy all defining relations of
$\St_n(T)$ and give a homomorphism
\[
 \Phi:\St_n(T)\longrightarrow P_0,
 \qquad x_{ij}(f)\longmapsto y_{ij}(f).
\]

Conversely, the assignment $g_{ij}(w)\mapsto x_{ij}(w)$ satisfies
Equations~\eqref{fp:finite-square}--\eqref{fp:finite-product}, and
hence defines a homomorphism $\Psi:P_0\to\St_n(T)$.  Recursive use
of Equation~\eqref{fp:steinberg-product} gives
$\Psi(y_{ij}(w))=x_{ij}(w)$ for every monomial $w$.  Additivity then
gives $\Psi(y_{ij}(f))=x_{ij}(f)$ for every $f\in T$.  It follows
that $\Psi\Phi$ fixes every defining generator of $\St_n(T)$.
The composition $\Phi\Psi$ fixes every defining generator
$g_{ij}(w)$ of $P_0$.  These two homomorphisms are inverse, proving
$P_0\cong\St_n(T)$.

\paragraph{Imposing the algebra ideal.}
Under this isomorphism, Equation~\eqref{fp:finite-algebra-relators}
kills $x_{ij}(f_t)$ for every $t$ and every root position.  Let $P$ be
the resulting quotient of $\St_n(T)$, and write $\widetilde x_{ij}(a)$
for the image in $P$ of $x_{ij}(a)$.  We verify directly that
\begin{equation}
 \widetilde x_{ij}(a)=1
 \quad\text{for every }a\in\mathfrak a\text{ and every }i\neq j.
 \label{fp:ideal-killed}
\end{equation}
For $c\in T$ and $h\notin\{i,j\}$, the product relation gives
\[
 \widetilde x_{ij}(cf_t)
 =[\widetilde x_{ih}(c),\widetilde x_{hj}(f_t)]=1.
\]
This conclusion holds in every root position.  Applying it to
position $(i,h)$ and using another product relation gives, for every
$d\in T$,
\[
 \widetilde x_{ij}(cf_td)
 =[\widetilde x_{ih}(cf_t),\widetilde x_{hj}(d)]=1.
\]
Every element of the two-sided ideal $\mathfrak a$ is a finite sum of
terms $cf_td$; additivity proves Equation~\eqref{fp:ideal-killed}.

The algebra quotient map $T\to R$ induces a surjective homomorphism
$\St_n(T)\to\St_n(R)$, and the specified relators lie in its
kernel.  It therefore induces a homomorphism $\alpha:P\to\St_n(R)$.
In the other direction, assign to $x_{ij}(r)$ the element
$\widetilde x_{ij}(a)$ for any lift $a\in T$ of $r\in R$.
This is independent of the lift: two lifts differ by an element of
$\mathfrak a$, whose root element is trivial by
Equation~\eqref{fp:ideal-killed} and additivity.  The Steinberg
relations hold for these elements by choosing lifts $a+b$ and $ab$
for sums and products.  Thus the assignment defines
$\beta:\St_n(R)\to P$.  The composition $\alpha\beta$ fixes
every $x_{ij}(r)$, and $\beta\alpha$ fixes every
$\widetilde x_{ij}(a)$, which generate $P$.  The two maps are inverse.

There are finitely many words of length at most two over $\mathcal A$,
finitely many indices, and finitely many algebra relators.  Every
$\xi_{ij}(w)$ appearing in the presentation is a finite word in the
specified generators.  Consequently all four displayed families give
finitely many ordinary group relators.  This proves the theorem.
\end{proof}

\subsection{Perfectness, infinitude, and matrix periodicity}

\begin{lemma}
\label{fp:perfectness}
For every unital associative $\F_2$-algebra $R$ and $n\geq3$, the
groups $\St_n(R)$ and $\E_n(R)$ are perfect.  If $R$ is infinite,
both groups are infinite.
\end{lemma}

\begin{proof}
For each $i\neq j$, choose $h\notin\{i,j\}$.  Every defining
generator of the Steinberg group is a commutator:
\[
 x_{ij}(a)=[x_{ih}(a),x_{hj}(1)].
\]
Hence the commutator subgroup contains all the generators, proving
perfectness.  Its quotient $\E_n(R)$ is perfect as well.  For a fixed
root position $(i,j)$, the matrices $e_{ij}(a)$ for distinct $a\in R$
are distinct, as is seen in their $(i,j)$ entries.  Thus an infinite
$R$ gives an infinite $\E_n(R)$, and the surjection
\eqref{fp:matrix-map} then gives an infinite $\St_n(R)$.
\end{proof}

\begin{lemma}[Periodicity of the matrix groups]
\label{fp:matrix-periodicity}
Suppose that $R=\F_2\oplus I$ as a vector space, where $I$ is a
two-sided ideal and the indicated copy of $\F_2$ contains the identity
of $R$.  Suppose that $M_d(I)$ is a nil algebra for every positive
integer $d$; that is, each individual matrix in $M_d(I)$ is nilpotent.
Then a square matrix over $R$ is invertible if and only if its reduction
modulo $I$ is invertible.  Moreover, every element of $\GL_n(R)$ has
finite order, and therefore $\E_n(R)$ is periodic.
\end{lemma}

\begin{proof}
An invertible matrix remains invertible under reduction.  Conversely,
let $A\in M_n(R)$ have invertible reduction $\overline A$ over
$\F_2$.  Regard $B=\overline A^{-1}$ as a matrix over the specified
copy of $\F_2$ in $R$; it is already invertible over $R$.  Then
$AB=\Id_n+U$ with $U\in M_n(I)$.  Choose $m\geq1$ such that
$U^m=0$.  In characteristic two,
\[
 (\Id_n+U)^{-1}=\Id_n+U+\cdots+U^{m-1}.
\]
Thus $AB$, and hence $A$, is invertible.

Now take $A\in\GL_n(R)$.  The finite group $\GL_n(\F_2)$ contains
its reduction, so $\overline A$ has some finite order $d\geq1$.
Consequently $A^d=\Id_n+V$ with $V\in M_n(I)$.  Choose $m\geq1$
with $V^m=0$, and choose $e\geq0$ with $2^e\geq m$.
Successive squaring in characteristic two gives
\[
 A^{d2^e}=(\Id_n+V)^{2^e}=\Id_n+V^{2^e}=\Id_n.
\]
This proves the periodicity assertions without requiring a common
nilpotence index or a common bound on element orders.
\end{proof}

By Equation~\eqref{fp:matrix-map}, if $\E_n(R)$ is periodic and each
element of $J_n(R)$ has finite order, then $\St_n(R)$ is periodic:
for $g\in\St_n(R)$, some positive power $g^d$ lies in $J_n(R)$,
and a further positive power is the identity.  The required control of
$J_n(R)$ is addressed in the subsequent sections.
